% ================= 4. DYNAMIC MODELING & CONTROL =================
\section{Dynamic Modeling \& Control}
\label{sec:modeling}

\subsection{Modeling Approach: Why 2D First}
\label{sec:2dfirst}

The full Cubli is an underactuated, nonlinear system that is unstable about all
three axes at once when balanced on a corner (Section~\ref{sec:objectives}).
Rather than commit directly to the fully coupled three-dimensional cube, this
project first develops and validates the complete modeling, identification, and
control workflow on a two-dimensional (2D) single-face prototype. The 2D rig
reduces the problem to a single rotational degree of freedom in a vertical plane,
which de-risks three things at the scale of one face before the full cube is
built: (i) the structure of the equations of motion and the control law, (ii) the
parameter-identification procedure that reconciles CAD-predicted properties with
the physical frame, and (iii) the jump-up mechanism. The dynamics validated on
the 2D model then generalize, axis by axis, to the 3D cube in
Section~\ref{sec:3dmodel}.

\subsection{2D Reaction-Wheel Pendulum}
\label{sec:2dmodel}

A two-dimensional (2D) prototype will be built at the scale of a single face of the cube, in order to develop and validate the control algorithms before the full cube is assembled. The prototype will consist of a 3D-printed plate matching the dimensions of one cube face, with a momentum wheel and motor mounted at its center and a braking mechanism at one corner. The plate will rest on a bearing at that corner, giving it a single rotational degree of freedom within a vertical plane.

% Suggested figure: free-body diagram of the 2D pendulum labeling
% theta_b, theta_w, l, l_b, and the pivot. This is the single highest-value
% addition for readability of the equations below.

\subsubsection{Equations of Motion}
\label{sec:2deom}

The nonlinear equations of motion governing the system will be
\begin{align}
\ddot{\theta}_b &= \frac{(m_b l_b + m_w l)\, g \sin\theta_b - T_m - C_b \dot{\theta}_b + C_w \dot{\theta}_w}{I_b + m_w l^2}, \\
\ddot{\theta}_w &= \frac{(I_b + I_w + m_w l^2)(T_m - C_w \dot{\theta}_w)}{I_w (I_b + m_w l^2)} \notag \\
&\quad - \frac{(m_b l_b + m_w l)\, g \sin\theta_b - C_b \dot{\theta}_b}{I_b + m_w l^2},
\end{align}
with the motor current-torque relation $T_m = K_m u$, where $\theta_b$ is the tilt angle of the pendulum body, $\theta_w$ the rotational displacement of the momentum wheel relative to the body, $m_b$ and $m_w$ the masses of the body and wheel, $I_b$ and $I_w$ their moments of inertia about the pivot and motor axis respectively, $l$ the distance from the motor axis to the pivot, $l_b$ the distance from the body's center of mass to the pivot, $g$ the gravitational acceleration, $T_m$ the motor torque, $C_b$ and $C_w$ the dynamic friction coefficients of the body and wheel, $K_m$ the motor torque constant of the selected motor, and $u$ the motor current input commanded through the motor driver.

\subsubsection{Equilibria \& Linearization}
\label{sec:2dequilibria}

% Identify the upright (edge-balance) equilibrium, linearize the equations of
% motion of Sec.~\ref{sec:2deom} about it, and present the state-space form used
% for the control design in Sec.~\ref{sec:control2d}.

\subsection{Parameter Identification \& the 2D Testbed}
\label{sec:paramid}

Since the pendulum body and wheel will be 3D-printed, their geometry will be fully defined in CAD prior to fabrication. This will allow $m_b$, $m_w$, $I_b$, $I_w$, and $l_b$ to be computed directly from the CAD model's mass-properties simulation, given the chosen print material, infill density, and the mass and placement of mounted hardware such as the motor driver, encoder, microcontroller, IMU, and wiring. These simulated values will be cross-checked experimentally once the prototype is fabricated: $l_b$ will be verified by freely hanging the body from different corners, and $I_b$ will be verified by fixing the wheel, hanging the assembly upside down, letting it swing, and fitting the identified model $(I_b+I_w+m_wl^2)\ddot\theta_b = -C_b\dot\theta_b + (m_bl_b+m_wl)g\sin\theta_b$ to the logged tilt data. The damping coefficients $C_b$ and $C_w$ cannot be obtained from simulation, since they depend on bearing friction, wiring drag, and motor cogging not captured in CAD, and will therefore be identified purely experimentally: $C_w$ (together with $I_w$, as a check against the simulated value) by driving the momentum wheel with current steps while the body is clamped and fitting $I_w\ddot\theta_w = K_mu - C_w\dot\theta_w$ to the logged wheel speed; $C_b$ from the swing-test data described above.

\begin{table}[h]
\centering
\caption{Prototype parameters and determination method.}
\label{tab:params}
\begin{tabular}{lll}
\toprule
Symbol & Description & Determination method \\
\midrule
$l$   & Motor axis to pivot distance & Design choice \\
$l_b$ & Body CoM to pivot distance   & Design choice; verified by hang test \\
$m_b$ & Pendulum body mass           & CAD mass-properties \\
$m_w$ & Momentum wheel mass          & CAD mass-properties \\
$I_b$ & Body moment of inertia       & CAD mass-properties; verified by swing test \\
$I_w$ & Wheel moment of inertia      & CAD mass-properties; verified by current-step fit \\
$C_b$ & Body damping coefficient     & Experimental (swing-test fit) \\
$C_w$ & Wheel damping coefficient    & Experimental (current-step fit) \\
$K_m$ & Motor torque constant        & Motor datasheet \\
\bottomrule
\end{tabular}
\end{table}

\subsection{Jump-Up Feasibility Analysis}
\label{sec:jumpup}

With $I_b$, $I_w$, $m_b$, $l_b$, $m_w$, and $l$ established from CAD simulation, the wheel angular velocity $\omega_w$ required for the jump-up maneuver will be estimated before fabrication. Assuming a perfectly inelastic collision between wheel and body, conservation of angular momentum during impact together with conservation of energy from the moment of impact to the edge-standing position gives
\begin{equation}
\omega_w^2 = (2-\sqrt2)\,\frac{I_w+I_b+m_wl^2}{I_w^2}\,(m_bl_b+m_wl)\,g.
\end{equation}
This simulation-derived value of $\omega_w$ will be used to verify, ahead of fabrication, that the selected motor can reach the required speed within the available bus voltage, and to size the wheel's inertia --- via its radius and any added rim mass --- so that the required speed remains within the motor's practical operating range without excessively reducing the recovery angle available for balancing. Once the prototype is built, the current-step and swing-test identification described above will be used to refine $\omega_w$ experimentally, together with an iterative trial-and-error adjustment to account for any deviation from the idealized inelastic-collision assumption (e.g., a non-zero coefficient of restitution or parameter uncertainty).

The braking mechanism will consist of a servo-driven barrier that is swung into a bolt head attached to the momentum wheel to stop it abruptly. The servo's response time will be a design constraint on the validity of the ``instantaneous stop'' assumption underlying the jump-up equation above; a slower stop than assumed will yield a lower effective post-impact body velocity than predicted, which will be compensated for through the same iterative tuning process rather than relying on the analytical estimate alone.
\newpage
\subsubsection{Torque-Limited Wheel-Inertia Target}
\label{sec:jumpup_target}

The idealized relation above returns the required wheel speed for a given
inertia, but it assumes an instantaneous, lossless momentum transfer. For
preliminary sizing we invert the problem---given the achievable motor speed,
\emph{what per-wheel inertia does the maneuver demand?}---using a momentum
budget that additionally penalizes the finite brake torque. The maneuver must
supply enough post-impact angular momentum to raise the cube's center of mass
from its resting contact to the balancing equilibrium, scaled by an energy
margin $\eta$. This defines an ideal per-wheel angular momentum
\begin{equation}
h_{w,\text{ideal}} = \sqrt{\eta}\;\frac{\sqrt{2 g\,\lambda\,\kappa}}{n}\;M L^{3/2},
\label{eq:hw_ideal}
\end{equation}
where $M$ and $L$ are the cube mass and edge length, $g$ is gravity, and
$\kappa$, $\lambda$, $n$ are dimensionless geometric factors set by the contact
case: the normalized moment of inertia about the pivot ($\kappa$), the relative
rise of the center of mass to the equilibrium ($\lambda$), and the
momentum-transfer factor of the contact ($n$). The corner case
($\kappa=\tfrac{11}{12}$, $\lambda=\tfrac{\sqrt3}{2}-\tfrac12$, $n=\sqrt2$)
governs the design as it is the more demanding of the two; the edge case
($\kappa=\tfrac{2}{3}$, $\lambda=\tfrac{1}{\sqrt2}-\tfrac12$, $n=1$) is less
strenuous.

A brake of finite torque $\tau_b$ cannot deliver the impulse instantaneously:
during the stop the cube must still overcome the gravitational tip-over torque
$\tau_g = MgL/2$, which sets a hard floor $\tau_b > \tau_g$ below which the cube
can never tip. The resulting loss is captured by a brake-amplification factor
\begin{equation}
\beta = \frac{\tau_b}{\tau_b - \tau_g},
\qquad
h_w = \beta\,h_{w,\text{ideal}},
\qquad
I_{w,\text{target}} = \frac{h_w}{\omega_{\max}},
\label{eq:Iw_target}
\end{equation}
where $\omega_{\max}$ is the maximum practical wheel speed. Table
\ref{tab:jumpup_budget} lists the resulting target for the current baseline
cube. This per-wheel inertia $I_{w,\text{target}}$ is the requirement that the
reaction-wheel design of Section~\ref{sec:rw_sizing} must meet.

\begin{table}[h]
\centering
\caption{Torque-limited wheel-inertia target (corner case, baseline cube).}
\label{tab:jumpup_budget}
\begin{tabular}{llr}
\toprule
Symbol & Description & Value \\
\midrule
$M$                & Cube mass                    & $1.45$ kg \\
$L$                & Cube edge length             & $180$ mm \\
$\eta$             & Energy margin                & $1.05$ \\
$\tau_b$           & Brake torque per wheel       & $5.0$ N\,m \\
$\tau_g$           & Gravitational tip-over floor & $1.28$ N\,m \\
$\beta$            & Brake-amplification factor   & $1.34$ \\
$\omega_{\max}$    & Max wheel speed              & $8400$ rpm ($880$ rad/s) \\
$h_w$              & Required angular momentum    & $0.277$ kg\,m$^2$/s \\
$I_{w,\text{target}}$ & \textbf{Target wheel inertia} & $\mathbf{3.15\times10^{-4}}$ \textbf{kg\,m}$^\mathbf{2}$ \\
\bottomrule
\end{tabular}
\end{table}

\subsection{Control Design on the 2D Model}
\label{sec:control2d}

\subsubsection{LQR Baseline Control}

The nonlinear dynamics will be linearized about the upright equilibrium $(\theta_b,\dot\theta_b,\dot\theta_w)=(0,0,0)$, giving
\begin{equation}
\dot x(t) = Ax(t) + Bu(t), \qquad x := (\theta_b,\dot\theta_b,\dot\theta_w),
\end{equation}
with $A, B$ built from the identified parameters $m_b, m_w, I_b, I_w, l, l_b, C_b, C_w, K_m$. Zero-order-hold discretization at a rate set relative to the open-loop unstable pole (with a safety margin) gives $x[k+1] = A_dx[k] + B_du[k]$.
\newpage
A discrete LQR gain $K_d = (K_{d1}, K_{d2}, K_{d3})$ will minimize
\begin{equation}
J(u) = \sum_k x^{\mathrm T}[k]Qx[k] + u^{\mathrm T}[k]Ru[k],
\end{equation}
giving the control law $u[k] = -K_d(\hat\theta_b[k], \hat{\dot\theta}_b[k], \hat{\dot\theta}_w[k])$ from IMU tilt/rate and encoder wheel-speed feedback.

To prevent a constant tilt-sensor offset from appearing as a nonzero steady-state wheel speed, a low-pass filter state $x_f[k+1] = (1-\alpha)x_f[k] + \alpha\hat\theta_b[k]$ will be subtracted from the tilt feedback:
\begin{equation}
u[k] = -K_d(\hat\theta_b[k]-x_f[k],\, \hat{\dot\theta}_b[k],\, \hat{\dot\theta}_w[k]).
\end{equation}
Any offset $d$ is then absorbed by $x_f$ at steady state. This fixed-gain design will serve as the baseline balancing controller.
\subsubsection{Gain Scheduling \& Mode Transitions}

The controller will switch from the open-loop jump-up sequence to the LQR balancing law once the estimated tilt enters a threshold window, $|\hat\theta_b|<\theta_{\text{th}}$, where the linearized model is expected to hold; a hysteresis band around $\theta_{\text{th}}$ will prevent chattering between modes if the body oscillates near the switching point. The baseline will use a single fixed gain $K_d$ computed at the upright equilibrium. If testing shows degraded recovery near the edges of the operating range — where the small-angle linearization becomes less accurate — gain scheduling will be added as an extension: additional gains $K_d^{(i)}$ computed at several linearization points across the tilt range, with the controller interpolating or switching between them based on $\hat\theta_b$, while keeping the same discrete structure and offset-correction filter.
% Handling of the transition between jump-up, capture, and balance regimes.

\subsubsection{Jump-Up Strategy}

Using the CAD-derived $I_b, I_w, m_b, l_b, m_w, l$, the wheel speed $\omega_w$ required for jump-up will be estimated before fabrication assuming a perfectly inelastic wheel-body collision, via conservation of angular momentum at impact and conservation of energy from impact to the edge-standing position:
\begin{equation}
\omega_w^2 = (2-\sqrt2)\,\frac{I_w+I_b+m_wl^2}{I_w^2}\,(m_bl_b+m_wl)\,g.
\end{equation}
This value will be checked against the selected motor's achievable speed at the available bus voltage, and used to size wheel inertia (radius, rim mass) so the required speed stays in range without excessively reducing the balancing recovery angle. After fabrication, $\omega_w$ will be refined experimentally with iterative trial-and-error tuning to account for deviations from the ideal inelastic-collision assumption. The braking barrier's servo response time will bound how well the "instantaneous stop" assumption holds; a slower stop will be compensated through the same tuning process.

% Control sequence that spins up the wheel and triggers the brake to execute the
% jump-up analyzed in Sec.~\ref{sec:jumpup}.

\subsection{Generalization to the 3D Cube}
\label{sec:3dmodel}

The dynamics, identification workflow, and control approach validated on the 2D
model extend to the full three-dimensional cube, where three orthogonal
reaction wheels must stabilize the body about its edge and corner equilibria.
\newpage
\subsubsection{Full-Body Equations of Motion}

Let $R \in SO(3)$ denote the orientation of the cube body relative to the world frame, $\omega_b \in \mathbb{R}^3$ the body's angular velocity expressed in the body frame, and $\omega_w = (\omega_{w,1},\omega_{w,2},\omega_{w,3})^{\mathrm T}$ the angular velocities of the three reaction wheels relative to the body, each about the axis $e_i$ normal to the face it is mounted on. The attitude kinematics are
\begin{equation}
\dot R = R\,[\omega_b]_\times,
\end{equation}
where $[\cdot]_\times$ is the skew-symmetric cross-product matrix.

The total angular momentum of the assembly about the pivot point, expressed in the body frame, generalizes the single-axis momentum term from the 2D model:
\begin{equation}
H = I\omega_b + \sum_{i=1}^{3} I_{w,i}\,\omega_{w,i}\,e_i,
\end{equation}
where $I$ is the $3\times3$ inertia tensor of the body-plus-wheels assembly about the pivot, generalizing $I_b + m_w l^2$. Euler's equation for a body under external torque gives
\begin{equation}
\dot H + \omega_b \times H = \tau_g - C_b\,\omega_b,
\end{equation}
where $\tau_g = m\,r_{cm}\times(R^{\mathrm T}\hat g)$ is the gravity torque about the pivot, $m$ the total assembly mass, $r_{cm}$ the body-frame vector from the pivot to the assembly's center of mass (the 3D generalization of $l_b$), $\hat g$ the world vertical unit vector, and $C_b$ a (generally diagonal) $3\times3$ damping matrix. Each wheel's dynamics generalize the single-wheel relation directly:
\begin{equation}
I_{w,i}\big(\dot\omega_{b,i} + \dot\omega_{w,i}\big) = T_{m,i} - C_{w,i}\,\omega_{w,i}, \qquad i=1,2,3,
\end{equation}
with $\omega_{b,i} = e_i^{\mathrm T}\omega_b$ and $T_{m,i} = K_{m,i}u_i$. Restricting motion to a single vertical plane and setting two of the three wheel torques to zero recovers the single-axis prototype model exactly, with $I_b$, $l_b$, $m_b$, $m_w$, $I_w$, $C_b$, $C_w$ corresponding to the relevant single-axis projections of $I$, $r_{cm}$, $m$, $I_{w,i}$, $C_{w,i}$.

\subsubsection{Axis Coupling}
In the 2D prototype, body and wheel dynamics were coupled through a single scalar reaction torque. In the 3D cube this coupling becomes fully multi-axis, arising from two sources. First, gyroscopic coupling enters through the $\omega_b \times H$ term: whenever the body has angular velocity about more than one axis, momentum stored in one wheel produces a reaction torque about the \emph{other} two axes even with no motor torque commanded there. This term is bilinear in $\omega_b$ and $H$ and vanishes only near equilibria where $\omega_b\approx 0$. Second, inertial coupling can arise if wheel placement or the body's own mass distribution introduces off-diagonal (product-of-inertia) terms in $I$. For a symmetric cube with wheels centered on each face along the face normals, $I$ is expected to be near-diagonal; any residual coupling from wiring, battery placement, or asymmetric hardware will be characterized using the same CAD-plus-experimental identification process developed for the 2D prototype, extended from a single scalar $I_b$ to a full inertia tensor. Because of these two coupling mechanisms, each wheel cannot be treated as an independent 1D subsystem, and the balancing controller must be designed on the coupled multi-axis model rather than as three decoupled copies of the 2D LQR baseline.
\newpage
% Explain how the single-axis 2D result generalizes and where cross-axis coupling
% appears as the contact geometry shrinks from an edge (line) to a corner (point).

\subsubsection{Edge vs.\ Corner Equilibria \& Linearization}

The cube's jump-up will proceed in the same two stages used in the original strategy: Flat-to-Edge, where a single wheel brakes to bring the cube onto one edge, followed by Edge-to-Corner, where a second wheel brakes to bring it onto a corner. Balancing on an edge constrains the body to rotate about a single fixed axis, so the dynamics reduce to the same single-axis form as the 2D prototype model, and the LQR baseline and offset-correction filter developed earlier apply directly, with $m_b$, $l_b$, $I_b$ replaced by their edge-balancing equivalents.

Balancing on a corner is genuinely three-dimensional: the body has three rotational degrees of freedom about the corner, and all three wheels contribute. Linearizing the full-body equations of motion about the corner-up equilibrium $(\theta,\omega_b,\omega_w)=(0,0,0)$, where $\theta\in\mathbb{R}^3$ parameterizes small orientation error from corner-up, gives a state-space model of the same form as the 2D case:
\begin{equation}
\dot x(t) = Ax(t) + Bu(t), \qquad x := (\theta,\omega_b,\omega_w) \in \mathbb{R}^9,
\end{equation}
with $A\in\mathbb{R}^{9\times9}$ and $B\in\mathbb{R}^{9\times3}$ built from the full inertia tensor $I$, the mass distribution $r_{cm}$, and the per-wheel parameters $I_{w,i}, C_{w,i}, K_{m,i}$. The discrete LQR design generalizes directly: $K_d$ becomes a $3\times9$ gain matrix rather than a $1\times3$ vector, with the same zero-order-hold discretization, quadratic cost, and offset-correction filtering reused, and $Q,R$ sized to weight all three tilt axes, body rates, and wheel speeds. The mode-transition logic will be extended with a third mode — Edge-to-Corner jump-up, triggered once edge balancing is stable — followed by a second threshold-based switch into the full 3D corner LQR once estimated orientation enters a neighborhood of the corner-up equilibrium.
% Linearize about the edge and corner equilibria; contrast stability structure
% and controllability with the 2D case.
